\chapter{Conclusions, Operational Recommendations \& Future Scope}
\label{chap:conclusion}

\section{Synthesis of Findings \& Key Research Contributions}
\label{sec:synthesis_of_findings}

This thesis has presented the conception, design, empirical benchmarking, and operational deployment of an end-to-end quantitative machine learning pipeline for **5-minute-ahead implied volatility forecasting ($\Delta \text{IV}_{t+5}$)** and **relative option valuation** on **NIFTY 50 weekly index options** traded on the National Stock Exchange of India (NSE). Conducted under the institutional mentorship of **Surabhi Investment Services Pvt. Ltd.**, the research resolves the persistence assumption paradox in intraday options market making and establishes a rigorous quantitative bridge between high-frequency machine learning and closed-form derivatives pricing.

The key empirical and methodological contributions of this work are summarized as follows:

\begin{enumerate}
    \item \textbf{Beating the High-Frequency Persistence Benchmark:}
    Treating implied volatility as a static random walk ($\widehat{\Delta \text{IV}}_{t+5} = 0$) is standard practice across intraday desks, yet leaves market makers exposed to adverse selection. By engineering multi-horizon autoregressive volatility lags (5m and 10m) and velocity indicators, our tuned Random Forest regressor achieves a **5.28\% reduction in MAE** (0.5214 vs.\ 0.5505) and a **7.72\% reduction in RMSE** (1.2188 vs.\ 1.3207) on out-of-sample test data, confirming Hypothesis 1 ($H_1$).

    \item \textbf{Demonstrating the Structural Inadequacy of Linear Models:}
    Ordinary Least Squares (OLS) linear regression suffered catastrophic failure, producing an out-of-sample MAE of **6.1720** and negative $R^2 = -0.2410$. This empirically demonstrates that the cross-sectional curvature of the implied volatility smile across moneyness cannot be captured by linear hyperplanes, confirming Hypothesis 2 ($H_2$) and establishing the necessity of non-linear ensembles.

    \item \textbf{Quantifying the Dominance of Volatility Velocity over Static Geometry:}
    Gini feature importance revealed that recent volatility momentum (\texttt{IV\_change\_prev\_5m} at 7.21\% and \texttt{IV\_change\_prev\_10m} at 6.74\%) and order-book liquidity spreads (\texttt{bid\_ask\_spread} at 3.52\%) account for the majority of predictive variance, far exceeding traditional static strike moneyness and maturity metrics.

    \item \textbf{Formulating Actionable Conviction Thresholds:}
    While raw directional predictions yield a noisy baseline accuracy of **53.20\%**, filtering for high-conviction forecast magnitudes ($|\widehat{\Delta \text{IV}}| \ge 0.10$) elevates directional hit-rate to **63.40\%** (and **64.47\%** at $\ge 0.25$), providing execution algorithms with a robust filter that eliminates 82\% of low-conviction noise churn, confirming Hypothesis 3 ($H_3$).

    \item \textbf{Operationalizing Relative Option Valuation:}
    The Relative IV Valuation Engine systematically maps predicted volatility spreads into three distinct economic regimes: **Fair/Neutral (77.92\%)**, **High IV/Rich (15.46\%)**, and **Low IV/Cheap (6.61\%)**, providing quantitative desks with immediate, actionable signals for credit spreads, debit spreads, and delta-neutral market making.

    \item \textbf{Multi-Factor Black-Scholes Stress Testing:}
    Comprehensive scenario analysis of near-ATM options demonstrated the dominance of Delta over Vega in short-dated contracts. A $+1\%$ spot move generates a $+188.83\%$ gain in option value, whereas a $+20\%$ relative volatility expansion produces only a $+21.33\%$ gain. A combined $-2\%$ market fall with a $+20\%$ volatility spike causes a $-94.84\%$ collapse in call value, proving that long vega cannot protect long calls in sharp directional market crashes.

    \item \textbf{Discovery of the Expiry-Day Structural Regime Breakdown:}
    Out-of-sample testing on the October 19, 2026 expiry afternoon ($T < 1.58\text{ hours}$) revealed a severe econometric breakdown: model MAE exploded from 0.52 to **17.42**, the persistence baseline outperformed machine learning, and directional accuracy collapsed to **19.09\%**. This proves Hypothesis 4 ($H_4$) and provides formal mathematical proof that continuous-time volatility models collapse under asymptotic gamma/theta singularities and vanishing vega in the final hours of weekly expiry.

    \item \textbf{Production Prototype Deployment:}
    The analytical system was operationalized via an automated 5-minute live prediction daemon interfacing with the Upstox API v2 and an institutional multi-page dark-themed Streamlit analytics dashboard.
\end{enumerate}

\section{Operational Recommendations for Institutional Trading Desks}
\label{sec:desk_recommendations}

Based on the empirical findings, the following operational protocols are recommended for proprietary trading and advisory desks, specifically at **Surabhi Investment Services Pvt. Ltd.**:

\subsection{1. Dynamic Quote Shading in Market Making}
Proprietary market-making algorithms quoting bid-ask spreads around mid-market should replace static persistence assumptions with model forecasts $\widehat{\text{IV}}_{t+5}$. When the model projects an upward volatility momentum shock ($\widehat{\Delta \text{IV}} > +0.10$), the automated quoter should shade its bid higher and widen its ask, capturing edge and avoiding toxic adverse selection by informed market orders.

\subsection{2. Mandatory Gating on Systematic Trade Execution}
Options desks executing directional volatility plays (long straddles, vega scalping) should enforce the **$\theta = 0.10$ IV Point Conviction Rule**. Restricting order entry exclusively to contracts where $|\widehat{\Delta \text{IV}}| \ge 0.10$ delivers a 63.40\% directional edge while eliminating over 80\% of brokerage and exchange transaction fee overhead.

\subsection{3. Relative Value Portfolio Construction}
Traders should systematically exploit the Relative IV Valuation screener:
\begin{itemize}
    \item When a strike is flagged as **Rich / Overvalued** ($\text{Spread} > +0.50$), initiate short credit spreads (e.g., sell the rich strike, buy a farther OTM wing) to capture implied volatility compression.
    \item When a strike is flagged as **Cheap / Undervalued** ($\text{Spread} < -0.50$), initiate long debit spreads or diagonal calendar spreads to capture prospective volatility expansion.
\end{itemize}

\subsection{4. Protocol for the Expiry-Day Regime Shift: The 2-Hour Kill Switch}
Desks must establish a strict risk rule: **automatically deactivate machine-learning volatility models exactly 2.0 hours prior to weekly contract settlement (13:30 IST on expiry day)**. As documented in Chapter \ref{chap:results}, non-expiry models suffer catastrophic covariate shift and numerical instability as $\tau \to 0$. From 13:30 to 15:30 IST, desks should transition exclusively to parity arbitrage, intrinsic value boundary tracking, and pure gamma-scalping execution rules.

\section{Limitations of the Current Study}
\label{sec:limitations}

While empirical results demonstrate robust predictive performance, several academic and operational limitations should be acknowledged:
\begin{enumerate}\setlength{\itemsep}{1pt}\setlength{\parskip}{0pt}
    \item \textbf{Temporal Sample Window:} The dataset covers weekly expiry cycles in October 2026. Long-term multi-year backtesting across bull markets, bear regimes, and high-volatility macro crises would provide broader longitudinal generalization.
    \item \textbf{Strike Geometry Restriction:} The study focused on liquid near-ATM options ($0.95 \le K/S \le 1.05$). Deep out-of-the-money wings (carrying extreme tail-risk information) were excluded to avoid quote illiquidity noise.
    \item \textbf{Sampling Frequency:} Regularized 5-minute snapshots were analyzed. While optimal for algorithmic intraday positioning and quote shading, sub-second HFT desks require tick-by-tick order-book depth to model millisecond queue dynamics.
    \item \textbf{Execution Slippage Modeling:} While microstructure spreads were integrated into feature engineering, live P\&L backtests did not simulate explicit execution queue positions or exchange connectivity latency.
\end{enumerate}

\section{Directions for Future Quantitative Research}
\label{sec:future_research}

The findings and infrastructure established in this thesis suggest several promising pathways for future academic and industry research:

\begin{enumerate}\setlength{\itemsep}{1pt}\setlength{\parskip}{0pt}
    \item \textbf{Deep Sequential Architectures \& Attention Mechanisms:} Future iterations could explore deep sequence models—such as Temporal Fusion Transformers (TFT) or Bidirectional LSTM networks with self-attention—trained on rolling sequences of order-book messages to capture long-range dependencies.
    \item \textbf{Explicit Markov-Switching Regime Models:} To resolve the expiry collapse, future research should develop a Markov-Switching Vector Autoregressive (MS-VAR) framework that models the continuous transition between the Normal Trading Regime ($T > 6\text{h}$) and the Hyper-Gamma Expiry Regime ($T < 2\text{h}$) as an endogenous latent state.
    \item \textbf{Cross-Asset Volatility Spillover Networks:} Expanding the feature space to include cross-asset signals—such as BankNIFTY implied volatility, India VIX, USD/INR currency option volatility, and overnight US S\&P 500 VIX movements—would capture macro risk contagion across correlated financial markets.
    \item \textbf{Reinforcement Learning for Automated Delta-Vega Hedging:} Rather than utilizing static heuristic trading rules, future work could train a Deep Q-Network (DQN) or Proximal Policy Optimization (PPO) agent that dynamically adjusts portfolio delta and vega hedges conditioned on live model predictions under realistic transaction friction.
\end{enumerate}
